\documentclass[11pt]{article}
\usepackage{amsmath, amssymb}
\usepackage{geometry}
\usepackage{booktabs}
\usepackage{hyperref}
\geometry{margin=1in}

\title{Advertiser Values and Externalities:\\Construction from Raw Tweet Data}
\date{}

\begin{document}
\maketitle

\section{Raw Data}

The dataset consists of $T = 11{,}853$ tweets (after filtering) collected as of $D = \text{February 2, 2025}$.
Each tweet $t$ carries the following raw fields:

\begin{center}
\begin{tabular}{lll}
\toprule
Symbol & Field & Description \\
\midrule
$I_t$ & \texttt{impression\_count} & Total number of times the tweet was shown \\
$A_t$ & \texttt{action\_count} & Total engagement actions (likes, retweets, replies, quotes, bookmarks) \\
$d_t$ & \texttt{days\_old} & Age in days: $(D - \text{created\_at}_t)$ \\
\bottomrule
\end{tabular}
\end{center}

\noindent
Age in months is defined as
\begin{equation}
    M_t = \frac{d_t}{30}.
\end{equation}

\section{Community Notes Externality Score}

\subsection{Note Classification}

Each note $n$ attached to tweet $t_n$ carries a binary classification $c_n$.
This is mapped to a signed score:
\begin{equation}
    s_n =
    \begin{cases}
        -1 & c_n = \texttt{MISINFORMED\_OR\_POTENTIALLY\_MISLEADING} \\
        +1 & c_n = \texttt{NOT\_MISLEADING}
    \end{cases}
\end{equation}

\subsection{Rating Aggregation}

Each note $n$ receives ratings from a set of raters $\mathcal{R}(n)$.
Each rating $r \in \mathcal{R}(n)$ has a helpfulness level mapped to a score:
\begin{equation}
    \phi_r =
    \begin{cases}
        -1 & \texttt{helpfulnessLevel} = \texttt{NOT\_HELPFUL} \\
         0 & \texttt{helpfulnessLevel} = \texttt{SOMEWHAT\_HELPFUL} \\
        +1 & \texttt{helpfulnessLevel} = \texttt{HELPFUL}
    \end{cases}
\end{equation}

The aggregate rating score for note $n$ is
\begin{equation}
    R_n = 1 + \sum_{r \in \mathcal{R}(n)} \phi_r,
\end{equation}
where the additive $1$ accounts for the original note author being treated as a rater
who agrees with their own note.

\subsection{Note and Tweet Externality Scores}

The signed externality contribution of note $n$ is
\begin{equation}
    X_n = s_n \cdot R_n.
\end{equation}

The raw externality score of tweet $t$ is the sum over all its notes
$\mathcal{N}(t)$:
\begin{equation}
    X_t^{\text{raw}} = \sum_{n \in \mathcal{N}(t)} X_n,
\end{equation}
with $X_t^{\text{raw}} = 0$ for tweets that received no notes.

\subsection{Normalization by Impressions}

To make the externality comparable across tweets with different reach, it is expressed
per 1{,}000 impressions and then normalized relative to the dataset mean.
The normalization constant is computed over all $T$ tweets before any filtering:
\begin{equation}
    \gamma
    = \frac{\bar{X}^{\text{raw}}}{\bar{I}/1000}
    = \frac{(1/T)\sum_t X_t^{\text{raw}}}{(1/T)\sum_t I_t / 1000}.
\end{equation}

For each tweet with $I_t > 0$, the normalized externality per 1{,}000 impressions is
\begin{equation}
    \hat{x}_t
    = \frac{X_t^{\text{raw}} / (I_t / 1000)}{\gamma}
    = \frac{X_t^{\text{raw}}}{\gamma \cdot (I_t / 1000)}.
    \label{eq:hat_x}
\end{equation}
A tweet with exactly average externality per impression has $\hat{x}_t = 1$.
Positive values indicate content deemed not misleading by the community;
negative values indicate content deemed misleading.

\noindent
\textbf{Outlier removal.} Tweets with more than 150 actions per 1{,}000 impressions
($A_t / (I_t / 1000) \geq 150$) are removed as outliers.

\section{Advertiser Value Score and Externality Score}

Both scores are converted to monthly rates using the tweet's impression velocity.

\subsection{Externality Score}

The externality score $e_t^{\text{score}}$ (units: normalized externality per month) is
\begin{equation}
    e_t^{\text{score}}
    = \hat{x}_t \cdot \frac{I_t / M_t}{1000}
    = \frac{X_t^{\text{raw}}}{\gamma \cdot (I_t/1000)} \cdot \frac{I_t}{1000 \cdot M_t}
    = \frac{X_t^{\text{raw}}}{\gamma \cdot M_t}.
    \label{eq:e_score}
\end{equation}
Note that $I_t$ cancels: the score depends only on the total raw note signal $X_t^{\text{raw}}$
and tweet age $M_t$, not on raw impression count.

\subsection{Advertiser Value Score}

The value score $v_t^{\text{score}}$ (units: actions per month) is
\begin{equation}
    v_t^{\text{score}}
    = \frac{A_t / (I_t / 1000)}{1} \cdot \frac{I_t / M_t}{1000}
    = \frac{A_t}{M_t}.
    \label{eq:v_score}
\end{equation}
Again $I_t$ cancels: the value score is simply the total engagement action count divided
by tweet age in months.

\section{Cost Parameterization}

The optimizer (\texttt{Collateralized\_Auction\_genetic\_script.py}) scales the raw scores
by two free parameters before running the auction:
\begin{equation}
    v_t = \zeta_v \cdot v_t^{\text{score}}, \qquad
    e_t = \zeta_e \cdot e_t^{\text{score}},
    \label{eq:scaling}
\end{equation}
where:
\begin{itemize}
    \item $\zeta_v$ (\texttt{action\_cost}) is the monetary value of one advertiser action;
          default $\zeta_v = 1$.
    \item $\zeta_e$ (\texttt{externality\_cost}) is the weight placed on one unit of
          Community Notes externality; varied across experiments as
          $\zeta_e \in \{0.001, 0.005, 0.0075, 0.01, 0.025, 0.05, 0.1\}$.
\end{itemize}

\section{Joint Normalization to $[-1,\,1]$}

The GA searches for $\tau$ in a normalized space to ensure numerical stability.
Both $v_t$ and $e_t$ are scaled jointly using the global minimum and maximum
across both columns:
\begin{equation}
    x_{\min} = \min\!\left(\min_t v_t,\; \min_t e_t\right), \qquad
    x_{\max} = \max\!\left(\max_t v_t,\; \max_t e_t\right).
\end{equation}
The normalized values are
\begin{equation}
    \tilde{v}_t = \frac{2(v_t - x_{\min})}{x_{\max} - x_{\min}} - 1, \qquad
    \tilde{e}_t = \frac{2(e_t - x_{\min})}{x_{\max} - x_{\min}} - 1,
    \label{eq:normalization}
\end{equation}
so that $\tilde{v}_t, \tilde{e}_t \in [-1,1]$ for all $t$ in the dataset.
Because the joint extrema are used, the relative scale between $v$ and $e$ is
preserved after normalization.

\section{Auction Mechanism}

\subsection{Advertiser Sampling}

For each of $K$ simulated auctions, $n$ tweets are drawn without replacement from the
full dataset. Each sampled tweet $i$ participates as an advertiser with pair
$(\tilde{v}_i, \tilde{e}_i)$ in normalized space.

\subsection{Threshold Function $\tau$}

The collateralized auction uses a polynomial threshold function of degree $d$:
\begin{equation}
    \tau(\tilde{e};\,\boldsymbol{\beta})
    = \sum_{k=0}^{d} \beta_k\, \tilde{e}^k,
    \qquad \boldsymbol{\beta} = (\beta_0, \ldots, \beta_d) \in \mathbb{R}^{d+1}.
\end{equation}
The polynomial degree $d \in \{1, 2, 3\}$ is a hyperparameter varied across experiments.

\subsection{Collateralized vs.\ VCG Auction}

Both auctions allocate $k$ slots to the highest-value bidders.

\begin{description}
    \item[VCG (counterfactual).] All $n$ advertisers bid their value $\tilde{v}_i$.
    The top-$k$ advertisers by $\tilde{v}_i$ win.

    \item[Collateralized.] Advertiser $i$ is admitted to the auction only if
    \begin{equation}
        \tilde{v}_i \;\geq\; \tau(\tilde{e}_i;\,\boldsymbol{\beta}).
    \end{equation}
    Among admitted advertisers, the top $k$ by $\tilde{v}_i$ win.
    If fewer than $k$ advertisers are admitted, all of them win.
\end{description}

\subsection{Welfare Objective}

For a set of $k$ winning advertisers $\mathcal{W}$, define the auction welfare as
\begin{equation}
    W(\mathcal{W}) = \sum_{i \in \mathcal{W}} \tilde{v}_i + \sum_{i \in \mathcal{W}} \tilde{e}_i.
\end{equation}
The fitness function used by the genetic algorithm is the mean collateralized welfare
averaged over all $K$ auction draws:
\begin{equation}
    F(\boldsymbol{\beta})
    = \frac{1}{K} \sum_{j=1}^{K} W\!\left(\mathcal{W}_j^{\text{coll}}(\boldsymbol{\beta})\right).
\end{equation}
The GA searches for $\boldsymbol{\beta}^* = \arg\max_{\boldsymbol{\beta}} F(\boldsymbol{\beta})$.

\subsection{Coefficient Recovery in Original Space}

After the GA returns $\boldsymbol{\beta}^*$ in normalized space, the optimal threshold is
converted back to the original $(e, v)$ coordinates using the inverse of the affine
transform in \eqref{eq:normalization}. This is done by
\texttt{NegOneOneScaler.poly\_from\_normalized}, which applies a change of variables
via Horner's rule to produce coefficients $\boldsymbol{\alpha}^*$ such that
\begin{equation}
    \tau^*(e) = \sum_{k=0}^{d} \alpha_k^*\, e^k
    \quad\Longleftrightarrow\quad
    \tau(\tilde{e};\,\boldsymbol{\beta}^*),
\end{equation}
where $e$ and $\tilde{e}$ are related by \eqref{eq:scaling} and \eqref{eq:normalization}.

\end{document}
